\chapter{Nonlinear Reinforced-Concrete 5-Story Building Example}
\label{ch:regular-5story-rc-nonlinear-building}

\section{Purpose and Scope}

This chapter documents the current \texttt{main.cpp} entry point, which now
hosts a materially nonlinear reinforced-concrete building example.  The model
is the direct continuation of the elastic mixed-building benchmark of
Chapter~\ref{ch:regular-5story-building}, but the frame members are now
represented by nonlinear fiber sections while the floor slabs remain elastic
MITC4 shells.

The example is intended to validate three architectural goals simultaneously:
\begin{itemize}
  \item one shared \texttt{Domain<3>} and one shared \texttt{DMPlex} for
        columns, beams, and slabs;
  \item one polymorphic structural element container, mixing beam and shell
        formulations in the same global solve;
  \item one constitutive upgrade path in which the frame material law is
        replaced by a nonlinear fiber discretisation without changing the
        geometric or topological assembly pipeline.
\end{itemize}

Equally important, the example is \emph{not} yet a fully nonlinear RC building
solver in the broadest sense.  The current scope is:
\begin{itemize}
  \item small-strain structural kinematics;
  \item nonlinear beam sections through fibers;
  \item elastic shell sections for slabs;
  \item static incremental equilibrium with PETSc SNES.
\end{itemize}

\section{Unit System}

The example switches deliberately to:
\[
\text{force} = \mathrm{MN},
\qquad
\text{length} = \mathrm{m},
\qquad
\text{stress} = \mathrm{MPa} = \mathrm{MN/m^2}.
\]

This choice is not cosmetic.  The uniaxial constitutive laws currently used in
the fiber sections, especially Kent--Park concrete, are calibrated and written
numerically in MPa.  Using \(\mathrm{MN}\)--\(\mathrm{m}\)--\(\mathrm{MPa}\)
preserves unit consistency while avoiding ad hoc rescalings inside the
constitutive model.  In this system:
\begin{itemize}
  \item \(30\ \mathrm{MPa}\) is entered numerically as \(30\),
  \item \(200\ \mathrm{kN}\) is entered as \(0.20\ \mathrm{MN}\),
  \item \(7.5\ \mathrm{kN/m^2}\) is entered as \(0.0075\ \mathrm{MN/m^2}\).
\end{itemize}

\section{Geometry and Topology}

\subsection{Plan and elevations}

The plan grid is unchanged from the elastic reference:
\[
x = \{0,\; 6,\; 12,\; 18\}\ \mathrm{m},
\qquad
y = \{0,\; 5,\; 10\}\ \mathrm{m},
\]
with storey height
\[
h = 3.20\ \mathrm{m},
\qquad
z_k = 3.20\,k,\quad k=0,\dots,5.
\]

This yields:
\begin{itemize}
  \item \(72\) shared building nodes,
  \item \(60\) columns,
  \item \(85\) beams,
  \item \(30\) shell panels,
  \item \(175\) geometric elements in total.
\end{itemize}

\subsection{Mixed structural topology}

All structural members are inserted into the same \texttt{Domain<3>}:
\begin{itemize}
  \item columns and beams as 1D line geometries,
  \item slabs as 2D quadrilateral geometries.
\end{itemize}

The resulting \texttt{DMPlex} remains non-interpolated, but its highest
topological dimension is now \(2\) because shell cells are present.  The
recent \texttt{DMPlex} hardening from Chapter~\ref{ch:dmplex-mixed-topology}
is what makes this mixed 1D/2D configuration natural rather than exceptional.

\section{Constitutive Design}

\subsection{Columns}

Columns are modelled as 3D Timoshenko beam elements with section:
\[
b_c = h_c = 0.45\ \mathrm{m}.
\]

The nonlinear section is built from:
\begin{itemize}
  \item unconfined cover concrete, modelled with Kent--Park;
  \item confined core concrete, also modelled with Kent--Park but with a
        confinement enhancement factor derived from an assumed volumetric tie
        ratio;
  \item eight longitudinal reinforcing bars of diameter
        \(\phi = 25\ \mathrm{mm}\), modelled with Menegotto--Pinto steel.
\end{itemize}

The material parameters are:
\[
f'_c = 35\ \mathrm{MPa},
\qquad
E_c = 1000\,f'_c = 35000\ \mathrm{MPa},
\qquad
\nu = 0.20,
\]
\[
E_s = 200000\ \mathrm{MPa},
\qquad
f_y = 420\ \mathrm{MPa},
\qquad
b_s = 0.01.
\]

The cover thickness to fiber centroid is taken as:
\[
c_c = 0.05\ \mathrm{m}.
\]

The column section is discretised into:
\begin{itemize}
  \item \(84\) concrete fibers,
  \item \(8\) steel fibers,
  \item total: \(92\) fibers per column integration point.
\end{itemize}

\subsection{Beams}

Beams are also 3D Timoshenko elements, with section:
\[
b_b = 0.30\ \mathrm{m},
\qquad
h_b = 0.60\ \mathrm{m}.
\]

The beam fiber section uses:
\begin{itemize}
  \item unconfined Kent--Park concrete over the full section;
  \item six longitudinal reinforcing bars of diameter
        \(\phi = 20\ \mathrm{mm}\), arranged as three top and three bottom
        bars;
  \item the same Menegotto--Pinto steel law as in the columns.
\end{itemize}

The beam concrete parameters are:
\[
f'_c = 30\ \mathrm{MPa},
\qquad
E_c = 30000\ \mathrm{MPa}.
\]

The beam section discretisation contains:
\begin{itemize}
  \item \(72\) concrete fibers,
  \item \(6\) steel fibers,
  \item total: \(78\) fibers per beam integration point.
\end{itemize}

\subsection{Slabs}

The slabs remain elastic MITC4 / Mindlin--Reissner shells with thickness:
\[
t = 0.15\ \mathrm{m}.
\]

The elastic shell properties are:
\[
E_s = 28000\ \mathrm{MPa},
\qquad
\nu_s = 0.20.
\]

This is a deliberate modelling boundary.  The current example is a nonlinear
\emph{frame} with an elastic slab diaphragm, not yet a fully fiberised RC
shell--frame interaction problem.

\section{Why Fibers Belong in the Constitutive Layer}

The crucial architectural decision is that the fiber discretisation is not a
visualisation gadget.  It is the constitutive relation itself:
\[
\texttt{FiberSection3D}
\;\rightarrow\;
\texttt{Material<TimoshenkoBeam3D>}
\;\rightarrow\;
\texttt{BeamElement}.
\]

That means the same reduction policy that maps generalized beam strain
\((\varepsilon_0,\kappa_y,\kappa_z,\gamma_y,\gamma_z,\theta')\) to section
response is already the one that will later be reusable for:
\begin{itemize}
  \item post-processing,
  \item section-force recovery,
  \item multiscale or scale-bridging initialisation,
  \item transfer of structural states to continuum models.
\end{itemize}

The visualisation layer only reads committed constitutive snapshots; it does
not define the physics.

\section{Section Kinematics and Fiber Stress Recovery}

For each beam material point, the generalized strain vector is:
\[
\mathbf{e} =
\begin{bmatrix}
\varepsilon_0 &
\kappa_y &
\kappa_z &
\gamma_y &
\gamma_z &
\theta'
\end{bmatrix}^{T}.
\]

The axial strain at a fiber located at local section coordinates \((y,z)\) is
reconstructed as:
\[
\varepsilon_f(y,z) = \varepsilon_0 - z\,\kappa_y + y\,\kappa_z.
\]

The section resultants then follow from numerical integration over the fibers:
\[
N   = \sum_i \sigma_i A_i,
\qquad
M_y = \sum_i \sigma_i(-z_i)A_i,
\qquad
M_z = \sum_i \sigma_i y_i A_i.
\]

The shear and torsional terms remain elastic:
\[
V_y = \kappa_y G A\,\gamma_y,
\qquad
V_z = \kappa_z G A\,\gamma_z,
\qquad
T   = GJ\,\theta'.
\]

This is why the current example is materially nonlinear but still efficient:
only the axial/flexural part of the section is fiber-integrated, while the
shear/torsion part remains compact and closed form.

\section{Global Solution Strategy}

The mixed model is solved with
\texttt{NonlinearAnalysis<..., StructuralPolicy>}, i.e. PETSc SNES on the same
polymorphic structural container already used by the linear example.  The load
vector is applied incrementally:
\[
\lambda_k = \frac{k}{12},
\qquad
k = 1,\dots,12,
\]
with automatic step bisection enabled up to depth \(6\), although the final
run converged without exhausting that budget.

The PETSc options used in the executable are:
\begin{itemize}
  \item \texttt{-snes\_type newtonls},
  \item \texttt{-snes\_linesearch\_type bt},
  \item \texttt{-snes\_rtol 1e-8},
  \item \texttt{-snes\_atol 1e-9},
  \item \texttt{-snes\_max\_it 40},
  \item \texttt{-ksp\_type preonly},
  \item \texttt{-pc\_type lu}.
\end{itemize}

\subsection{Analysis-layer hardening}

In order for this example to compile and solve through the structural
type-erased path, \texttt{NLAnalysis} was generalised.  The solver now uses two
assembly routes:
\begin{itemize}
  \item a fast local-vector route for continuum elements exposing explicit
        element DOF extraction and local tangent routines;
  \item a generic \texttt{FiniteElement} route for structural type-erased
        elements, assembling through \texttt{compute\_internal\_forces()} and
        \texttt{inject\_tangent\_stiffness()}.
\end{itemize}

This keeps the continuum hot path intact while making the nonlinear solver
architecturally consistent with the structural abstraction already present in
the codebase.

\section{Loads and Boundary Conditions}

The support conditions are unchanged from the baseline model:
all nodes on \(z=0\) are fully fixed.

The slabs carry a uniform gravity pressure:
\[
q_g = -0.0075\ \mathrm{MN/m^2}.
\]

The lateral load pattern is triangular and applied at the corner
\((18,10,z_k)\):
\[
F_x^{(k)} = 0.20\,k \quad \mathrm{MN},
\qquad
k = 1,\dots,5.
\]

This choice keeps the model response moderate while making sway and beam
flexure clearly visible.

\section{Results of the Executed Example}

The example was executed successfully in the current workspace.  The final
incremental solve converged at \(\lambda=1.0\) with \(4\) SNES iterations in
the last step.

The reported response values are:
\[
\max \|\mathbf{u}\| = 0.003767\ \mathrm{m},
\qquad
\max |u_x| = 0.003314\ \mathrm{m},
\qquad
\max |u_z| = 0.001244\ \mathrm{m},
\]
and at the roof corner:
\[
u_x = 0.003314\ \mathrm{m},
\qquad
u_y = -0.001473\ \mathrm{m},
\qquad
u_z = -0.001019\ \mathrm{m}.
\]

The roof drift ratio is:
\[
\frac{u_x^\text{roof}}{H} = 2.07\times 10^{-4}.
\]

This is intentionally a mild nonlinear response, not a collapse or pushover
limit-state study.  The objective here is robustness and architectural
traceability, not maximum inelastic excursion.

\section{VTK Output}

The example writes:
\begin{verbatim}
regular_5story_rc_nonlinear_building_structural.vtm
\end{verbatim}
plus the corresponding block sidecars:
\begin{itemize}
  \item \texttt{\_axis.vtp}
  \item \texttt{\_surface.vtp}
  \item \texttt{\_sections.vtp}
  \item \texttt{\_material\_sites.vtp}
  \item \texttt{\_fibers.vtp}
\end{itemize}

For nonlinear RC interpretation, the most important block is:
\begin{itemize}
  \item \texttt{fibers}: it contains the actual fiber-level
        \texttt{fiber\_strain\_xx} and \texttt{fiber\_stress\_xx} values
        reconstructed from the committed constitutive state.
\end{itemize}

The surface block remains useful for displacement visualisation and geometric
context, but the truly nonlinear stress information lives at the fiber level.

\section{Regression and Verification}

Alongside the full building example, a dedicated regression test was added:
\begin{verbatim}
tests/test_rc_fiber_frame_nonlinear.cpp
\end{verbatim}

It solves a small RC fiber cantilever with the same
\texttt{BeamElement + FiberSection + NonlinearAnalysis} stack and checks:
\begin{itemize}
  \item convergence of the incremental SNES solve,
  \item finite tip displacement,
  \item physically correct sign of the vertical response.
\end{itemize}

During that work a real lifecycle bug was found and fixed in the test harness:
PETSc-owning objects must be destroyed before \texttt{PetscFinalize()}.

\section{Current Limits and Next Steps}

The example is already useful as a production-quality integration case, but
some limits remain explicit:
\begin{itemize}
  \item the slab is still elastic;
  \item the structural kinematics are still small-strain;
  \item beam surface stresses in the generic surface carrier are only a
        kinematic proxy when no closed-form beam snapshot exists, whereas the
        fiber block contains the actual nonlinear stress information;
  \item confined/unconfined concrete zoning is still schematic rather than
        design-code calibrated.
\end{itemize}

The next natural architectural step is not in visualisation, but in
reconstruction: extend the same reduction-policy logic toward consistent
scale-bridging and, eventually, transfer of committed structural states to a
continuum discretisation.
